\begin{proof}

\begin{enumerate}[label={$\langle1\rangle$\arabic*.},leftmargin=*]
\item The component $B$ has at least three generic normalization
branches, hence transverse multiplicity at least three.
Along the smooth $A$, choose coordinates $A=\{x=y=0\}$,
with $z$ on $A$. The actual equation belongs to $(x,y)^3$;
write $q_3(x,y;z)$ for its normal cubic.

\item A local branch of $B$ has first nonzero normal displacement
$(a,b)t^q$, with $(a,b)\ne0$. All ambient Hessian entries
of the actual equation vanish identically along $B$, by
its generic multiplicity. The first normal-order term of
its normal Hessian is
\[
       t^q\operatorname{Hess}_{x,y}
                    (q_3(\,\cdot\,,\,\cdot\,;0))(a,b).
\]
Terms of higher normal order and terms from varying $z$
have larger order in $t$. This matrix is therefore zero.
For a homogeneous binary cubic $Q$, a nonzero vector
annihilating this Hessian makes the directional derivative
of $Q$ the zero quadratic. In coordinates with that vector
as first axis, $Q$ is a cube of the other variable, or zero.
P4R.19 excludes the nonzero cube. Thus the cubic vanishes
at every $A$--$B$ junction, including tangent junctions.

\paragraph{The global spectral equality.}
Let $p,q$ be the two distinct junctions. The common cubic
zero divisor has degree $b\ge2$. On $\mathbb P(N_{A/Y}^*)$
its residual spectral curve is $\Gamma_H\cup\Gamma_R$,
of class $3\xi-bF$. The bundle has degree two, $\Gamma_R$
is a section, and $H$ normalizes $\Gamma_H$. Adjunction
and the exact union sequence give
\[
       g(H)+\delta(\Gamma_H)+\Gamma_H\cdot\Gamma_R=5-2b.
\]
More precisely, the displayed formula is an equality of
integers with $\delta(\Gamma_H)$ the total normalization
length. Since $g(H)\ge1$, it forces
\[
 b=2,\qquad g(H)=1,\qquad \delta(\Gamma_H)=0,
                  \qquad\Gamma_H\cdot\Gamma_R=0.          \tag{4.57a}
\]
The two common cubic zeros are simple, the double spectral
cover is smooth, and it is disjoint from the third section.

\paragraph{The graph support at a junction.}
Let $c_B$ count the global conductor components over $B$.
The graph budget is $c_B+1$. At either junction the
contribution is at least $\min(r_A,r'_B)$, where
$r_A\ge2$ and $r'_B\ge c_B$ count normalized points
over its two image branches.

\item If $c_B=2$, the two contributions total at least four,
exceeding the budget three. If $c_B=1$, a junction with
at least two $B$ preimages costs at least two, and the
other at least one, exceeding the budget two. Thus both
have one $B$ preimage. The argument of P4.56 gives the
path $H$--$V$--$R$, with one edge at each junction.
At the $R$--$V$ junction the local branch carries all
nearby generic $B$ sheets of $V$, regardless of its
finite degree. Its intersection with the $H$ branch or
branches is a length-one or length-two thickening of
smooth $A$. P4.56 makes it smooth, contradicting its
intrinsic nonnormality from the at least three $B$ sheets.

\item Therefore $c_B\ge3$. If $A$ belongs to every branch
support at a junction, that point costs at least
$r'_B\ge c_B$, and the other costs at least two.
This exceeds $c_B+1$. The two-element support argument
consequently forces $B$ to belong to every local
hypersurface branch at each junction. Branches carrying
only $B$ have been retained in this conclusion.

\paragraph{Choosing distinct normalization points.}
The ruling descends to $r_N:N\to H$. Its restriction
to $H$ is the identity, and $R$ maps to a fixed point
$b_0\in H$. If a point of $H$ and a point of $R$ have
the same image in $N$, that $H$ point is $b_0$.
This includes identifications through a contracted
exceptional chain. All such coincidences therefore lie
over the single value $a_0\in A$ which is the image of
$b_0$. Choose a junction different from $a_0$, and call
it $p$. Its $H$ preimages and its $R$ preimage are
distinct normalization points. Different $H$ points
are already distinct because $H$ is embedded.

\item The actual local equation factors into the corresponding
irreducible hypersurface factors and any factors carrying
only $B$. For an $A$-containing factor, let $b_i\ge0$
be the common $z$-order of its first nonzero normal term.
For a factor not containing $A$, let $c_j\ge1$ be the
order of its nonzero restriction to $A$ at $p$.
Gauss's lemma over $\mathbb C\{z\}$ gives
\[
                         \sum_i b_i+\sum_j c_j=1.         \tag{4.57b}
\]
Indeed primitive homogeneous polynomials have nonzero
product modulo $z$, so their common coefficient orders
add exactly. The right side is the simple cubic zero
of (4.57a).

\item If $H$ is unramified here, there are three
$A$-containing factors of normal order one. At least
two have nonzero linear normal terms, by (4.57b).
Their directions differ by (4.57a). The implicit-function
theorem makes their intersection exactly $A$, whereas
both must contain $B$. This is a contradiction.
Thus $H$ is ramified at $p$.

\item There are now two $A$-containing factors, denoted by
$F_H$ and $F_R$, of normal orders two and one. Their
residual spectral directions differ. Equality (4.57b)
has precisely the following three possibilities.

\paragraph{A factor carrying only $B$ uses the zero.}
Here $b_H=b_R=0$ and one such factor restricts to $A$
with order one. The branch $F_R$ is smooth. In coordinates
$X=F_R$, $A=\{X=Y=0\}$, restriction of $F_H$ to $X=0$
is $Y^2$ times a unit, since its residual direction
differs from the $R$ direction. Thus these two branches
intersect only along $A$. They both contain $B$, which
is impossible.

\paragraph{The quadratic factor uses the zero.}
Here $b_H=1$, $b_R=0$, and there is no other factor.
The branch $F_R$ is again smooth, and in its coordinates
exact division gives
\[
                         F_H|_{X=0}=Y^2h(Y,z).
\]
The value $h(0,z)$ has a simple zero: the residual
quadratic is nonzero on the $R$ direction, so its
common scalar factor survives with order one on this
restriction. Thus $h_z(0,0)\ne0$. Its zero curve is the
extra image branch $B$, with intersection multiplicity
one between $F_H$ and $F_R$.

\item The smooth $F_R$ branch supplies one generic $B$
normalization branch. Each branch supplied by $F_H$
has positive intersection multiplicity with it. There
can therefore be at most one such branch. The full
$B$ germ would have at most two branches, contradicting
P4.56.

\paragraph{The linear factor uses the zero.}
Finally $b_R=1$, $b_H=0$, with no other factor.
The nonzero quadratic form of $F_H$ has a double root.
Preparation and depression preserving $A$ give
\[
                   F_H=\text{unit}\,[Y^2+X^2k(X,z)].
\]
Its spectral curve is $u^2+k(0,z)=0$. By (4.57a) this
is smooth and its double cover is ramified at $p$.
Hence $k_z(0,0)\ne0$. The exact finite map
\[
 \{u^2+k(X,z)=0\}\longrightarrow\{F_H=0\},
                         \qquad Y=Xu
\]
has smooth source, is birational, and is an isomorphism
off $A$. It is the normalization. Thus $F_H$ supplies
one generic $B$ branch.

\item For $F_R$, the simple common zero of its linear normal
term makes its quadratic Taylor part
\[
                       zL(X,Y)+Q_2(X,Y),\qquad L\ne0.
\]
There is no $z^2$ term because $F_R$ vanishes on $A$.
Its Hessian has rank at least two. Rank three gives an
isolated critical point. In rank two the holomorphic
Morse lemma with a parameter gives the exact equation
$U^2+V^2+h(W)$. A nonisolated critical curve forces
$h'\equiv0$, hence $h=0$ because the factor vanishes
at the point. The resulting factorization
$(U+iV)(U-iV)$ contradicts local irreducibility.
Therefore $F_R$ has at most isolated singularities and
also supplies one generic $B$ branch. Again the full
branch number is two, contradicting P4.56.

\item All allocations in (4.57b) have been excluded, proving
the assertion. The exact genus budget and nonzero
Hessian minor are recomputed in
\texttt{code/part4/d4\_two\_junction.py}; its saved
output reports PASS.

\end{enumerate}
\end{proof}
